\section{Decision timing and telemetry construction}
The common forecast uses upstream observations through index $t-6$ and predicts active power at $t+6$. All timestamps mark the end of a ten-minute interval. The reserve learner receives the common point forecast, sine and cosine of time of day, and normalized turbine identity. The delayed background occupies 36 fixed event-time slots ending at $t$; its six most recent slots are masked. The last available slot is therefore $t-6$. A missing slot remains missing; the window is not shifted to find an earlier observation.

For the 134-turbine comparison, the delayed background includes 36 slots each of wind speed, active power and reactive power. Six summaries are derived from the last six arrived active-power observations: mean, population standard deviation, ordinary least-squares slope, one-step change, six-step change and forecast minus last power. A missing required element makes its summary unavailable. The background therefore has four common columns, 108 history columns and six summaries, for 118 raw columns. The single current-power and current-pitch updates each add one column. The augmented-power update adds current power and two differences, for 121 raw columns. The evaluation population requires a finite future target, current active power and current pitch, which yields 671,538 matched intervals per fit over 35 days.

For Hill of Towie, the pitch-withheld family contains Common and the 21 nonempty combinations of wind speed, active power and reactive power at history lengths 1, 6 and 36. P1 adds current active power. WPQ36 uses 36 observations of all three channels. Four shared forecast and time/turbine columns remain in both configurations; the resource tuple counts additional channels, learned predictor calls and raw reserve input width.

\section{Source data and chronological evaluation}
SDWPF contains 134 turbine series sampled at ten-minute resolution \cite{zhou2024sdwpfdata}. Days 1--90 train the point forecaster. Days 136--175 train reserve rules, days 176--192 calibrate residual offsets, days 193--210 provide historical assessment, and days 211--245 form the evaluation window. Five matched fits use seeds 201--205 and 360,000 training rows per fit. The evaluation population includes 671,538 valid turbine intervals over 35 shared weather days.

The SDWPF window is a development evaluation. The cross-year study uses separately fitted Hill of Towie models. The two studies are not pooled into a common effect estimate.

Hill of Towie contains 21 turbines with UTC period-end records \cite{clerc2026towie}. The common point forecaster is fitted on January--June 2023 using up to 120,000 observations. Reserve rules are fitted on July--December 2023 using 360,000 matched observations per seed. Offset calibration uses January--June 2024, and historical selection uses July--December 2024. The 2025 evaluation contains 1,059,248 finite future active-power targets across 365 days. The January--April 2026 evaluation contains 361,326 targets across 120 days. Zero and negative power values are retained. The same future-target mask is applied to all compared rules.

The 2026 records include the documented Dynamic Yaw control campaign. We retain all prespecified turbines and times, so this window evaluates the fixed reserve rules across a change in operating distribution. It is not an estimate of the causal effect of the control campaign.

Hill of Towie wind speed is taken from \path{wtc_AcWindSp_mean}. Active and reactive power use \path{wtc_ActPower_mean} and \path{wtc_ReactPwr_mean}, respectively. Negative wind speed is marked unavailable; finite negative active and reactive power are retained. The future-target mask requires finite active power, irrespective of the reserve input configuration. For the matched SDWPF update, eligibility additionally requires finite current power and mean pitch. Missing time--turbine coordinates remain unavailable. No operating-state or control-campaign exclusion is applied.

\section{Estimation and calibration}
The shared point forecaster is a squared-error HGB fitted to the 36-observation wind-speed and active-power histories. At Hill of Towie it uses up to 120,000 historical rows and seed 20261018. Its fitted target is twelve steps after the last upstream observation. Training labels must mature before the first publication used for reserve fitting. The same maturity rule separates reserve fitting, calibration and historical selection. The point forecaster is held fixed across the reserve-input comparisons.

HGB uses quantile loss at 0.90, 100 iterations, at most 15 leaves, minimum leaf size 50, learning rate 0.08 and L2 regularization 1. MLP uses hidden widths 64 and 32 with ReLU, AdamW with learning rate 0.001 and weight decay 0.0001, 80 epochs and batch size 2048. The five fixed seeds are 201--205. HGB replaces unavailable values with zero and adds missingness indicators. MLP scaling uses finite training means and population scales, followed by zero imputation and missingness indicators. Preprocessing is fitted using training data only.

HGB early stopping is disabled. Its indicator count depends on historical missingness; the MLP retains one indicator for every raw input. The MLP target scale is the fitting-target standard deviation, floored at one, and its output bias uses the fitting-target 0.90 quantile. Training uses deterministic CPU settings. Raw widths in Table S1 precede these transformations and do not imply equal transformed widths or computation across estimators.

Both estimators learn the 0.90 shortfall quantile. Nonnegative raw reserve predictions are calibrated by a signed offset, equal to the higher empirical quantile of shortfall minus raw prediction on the designated historical calibration period. The registered offset grid is $q\in\{0.90,0.925,0.95,0.975,0.99\}$. Thus, $q$ identifies an offset level rather than a separately trained quantile model. The main fixed-reference comparison uses $q=0.95$; the secondary finite-family analysis uses $q=0.925$.

Historical offset selection uses a single level shared by the five fits of each recipe. Eligible levels satisfy the historical 10\% event-rate bound for every fit; selection minimizes their five-fit mean cost. The richer reference is chosen among historically risk-qualified WPQ36 recipes. Subsequent windows retain the selected input configuration, estimator and offset. The persistence comparator requests $\max(F-P^{\rm last},0)$ before applying the historical offset, using zero for an unavailable last power value. Its resource tuple is $(1,0,2)$, counting only the forecast and last power.

\section{Qualification and block inference}
For a candidate and reference, the relative mean-cost difference is $D_C=\bar C_{\mathrm{candidate}}/\bar C_{\mathrm{reference}}-1$. Qualification requires the simultaneous upper bound on $D_C$ to be no greater than 0.02. It also requires the event-rate upper bound for each of the five instances of both rules to be no greater than 0.10. The reference must have passed historical risk selection.

Intervals use 20,000 circular samples of contiguous UTC weather-day blocks. The primary block length is three days, and one-day and seven-day blocks are sensitivity analyses. The same sampled days are applied to both policies and all fitted instances. Daily numerators and valid-record counts are resampled together before pooled rates and costs are computed. Fits and turbines remain fixed. A maximum standardized deviation produces an approximate simultaneous 95\% envelope over each declared family \cite{kunsch1989jackknife}. The matched measurement study uses bootstrap seed 20261030. The Hill of Towie fixed-reference analyses use seed 20261029, with a separate 14-coordinate primary family in each year. Each primary family contains ten fit-specific event rates, two mean costs and two directed cost ratios. The measurement-study family additionally contains fit-specific shortfall-magnitude contrasts. Simultaneity applies within each declared family and block length. This inference is conditional on the fixed fits and assumes the observed weather-day sequence is informative about the evaluation window.

\section{Finite-family minimum identification}
The primary $q=0.95$ assessment compares one historically selected rule against a fixed reference. The secondary 2025 analysis evaluates all registered pitch-withheld combinations, estimators and offset levels. A policy is sufficient when its own event-rate upper bounds satisfy the cap and its paired cost-ratio upper bounds satisfy the margin against all possibly risk-feasible competitors. A smaller candidate is excluded only if an event-rate lower bound exceeds the cap or a paired cost lower bound exceeds the margin relative to a risk-qualified witness. Unresolved candidates are not counted as excluded.

At $q=0.925$ and the 2\% cost margin, the secondary analysis identifies P1/HGB with resource tuple $(1,1,5)$. It excludes 15 strictly smaller registered candidates, two by risk and 13 by cost. The primary $q=0.95$ comparison does not belong to that full-family sufficient set. The finite-family assessment was not repeated in 2026.
